% TOP_PROBLEM: {"id": 4800019, "title": "Multiple ergodic averages — Problem 19", "category_id": 11, "category": {"id": 11, "name": "dynamical_systems", "display_name": "Dynamical Systems", "description": "Problems about long-term behavior of deterministic systems, Hamiltonian dynamics, and periodic orbits.", "slug": "dynamical-systems", "order_index": 11, "created_at": "2026-07-31T15:26:25.671Z"}, "status": "partially_solved", "problem_number": "AMR-047-0019", "set_id": 13}
% FIRST_POSTED: 2026-10-06
% ATTEMPT_STATUS: unresolved
% OPENAI_CLAIM: {"schema_version":1,"source_repo":"https://github.com/openai/math","source_commit":"adc7f1241b42e322a6451854ab7e4b4c146bf78a","release_date":"2026-10-06","match":"partial","resolution":"partial","independently_reviewed":false,"second_pass":"n/a","families":[{"family":"154","title":"Pointwise multiple ergodic averages for mixing transformations","subject":"Dynamical systems and ergodic theory","manuscripts":[{"title":"Pointwise Multiple Ergodic Averages for Mixing Transformations","pdf_path":"preprints/Pointwise-Multiple-Ergodic-Averages-for-Mixing-Transformations-October-4-2026/multiple-ergodic-averages.pdf","readme_path":"preprints/Pointwise-Multiple-Ergodic-Averages-for-Mixing-Transformations-October-4-2026/README.md","date":"2026-10-04","theorem_ref":"thm:main (sections/introduction.tex:20)","statement_sources":[{"path":"preprints/Pointwise-Multiple-Ergodic-Averages-for-Mixing-Transformations-October-4-2026/build/sections/introduction.tex","line":20,"label":"thm:main"}]},{"title":"Pointwise convergence of fourfold ergodic averages for mixing transformations","pdf_path":"preprints/Pointwise-convergence-of-fourfold-ergodic-averages-for-mixing-transformations-October-4-2026/fourfold-ergodic-averages.pdf","readme_path":"preprints/Pointwise-convergence-of-fourfold-ergodic-averages-for-mixing-transformations-October-4-2026/README.md","date":"2026-10-04","theorem_ref":"thm:main (sections/01-introduction.tex:21)","statement_sources":[{"path":"preprints/Pointwise-convergence-of-fourfold-ergodic-averages-for-mixing-transformations-October-4-2026/build/sections/01-introduction.tex","line":21,"label":"thm:main"}]},{"title":"Pointwise convergence of triple ergodic averages for mixing transformations","pdf_path":"preprints/Pointwise-convergence-of-triple-ergodic-averages-for-mixing-transformations-October-4-2026/pointwise-triple-ergodic-averages-mixing-transformations.pdf","readme_path":"preprints/Pointwise-convergence-of-triple-ergodic-averages-for-mixing-transformations-October-4-2026/README.md","date":"2026-10-04","theorem_ref":"thm:main (paper/sections/introduction.tex:21)","statement_sources":[{"path":"preprints/Pointwise-convergence-of-triple-ergodic-averages-for-mixing-transformations-October-4-2026/build/paper/sections/introduction.tex","line":21,"label":"thm:main"}]},{"title":"Triple ergodic averages with distinct integer slopes","pdf_path":"preprints/Triple-ergodic-averages-with-distinct-integer-slopes-October-4-2026/triple-ergodic-distinct-slopes.pdf","readme_path":"preprints/Triple-ergodic-averages-with-distinct-integer-slopes-October-4-2026/README.md","date":"2026-10-04","theorem_ref":"thm:main (sections/introduction.tex:19)","statement_sources":[{"path":"preprints/Triple-ergodic-averages-with-distinct-integer-slopes-October-4-2026/build/sections/introduction.tex","line":19,"label":"thm:main"}]}],"lean":null,"reasoning_trace":null}]}
\documentclass[11pt]{article}
\usepackage[margin=25mm]{geometry}
\usepackage{fontspec}
\setmainfont{Latin Modern Roman}
\usepackage{amsmath,amssymb,mathtools}
\usepackage{unicode-math}
\setmathfont{Latin Modern Math}
\usepackage{enumitem}
\usepackage{xurl,hyperref,bookmark}
\hypersetup{colorlinks=true,urlcolor=blue}
\setcounter{secnumdepth}{0}
\setlength{\parindent}{0pt}
\setlength{\parskip}{5pt}
\setlength{\emergencystretch}{3em}
\begin{document}
\section*{Multiple ergodic averages — Problem 19}

\subsection{OpenAI's claim}
Manuscript 1 (family 154, 2026-10-04): For every invertible mixing transformation, consecutive multiple ergodic averages of any fixed finite length are claimed to have an almost-everywhere limit equal to the product of the input integrals, without a mixing rate; Manuscript 2 (family 154, 2026-10-04): Fourfold averages at times $n,2n,3n,4n$ are claimed to converge pointwise to the product of integrals for every invertible mixing system and bounded inputs on arbitrary probability spaces. Manuscript 3 (family 154, 2026-10-04): For any invertible mixing system, the threefold averages at $n,2n,3n$ are claimed to converge almost everywhere along all positive averaging lengths to the product of integrals; Manuscript 4 (family 154, 2026-10-04): Triple averages with any three distinct nonzero integer slopes, including negative ones, are claimed to converge almost everywhere to the product of integrals for mixing transformations.

\subsection{Mathematical statement of the claimed result}
\paragraph{Mathematical claim 1.} Let $(X,\mathcal F,\mu)$ be any probability space with an invertible bimeasurable measure-preserving transformation $T$. Assume mixing, meaning $\mu(A\cap T^{-m}B)\to\mu(A)\mu(B)$ as $|m|\to\infty$ for all measurable $A,B$. OpenAI claims, for every $n\ge2$ and every fixed bounded measurable $f_1,\ldots,f_n$,
\[
\frac1N\sum_{k=1}^N\prod_{j=1}^n f_j(T^{jk}x)\longrightarrow\prod_{j=1}^n\int_X f_j\,d\mu
\quad\text{for almost every }x.
\]
The limit uses all positive integers $N$, requires no mixing rate or standardness of the probability space, and its null exceptional set may depend on the functions.

\paragraph{Mathematical claim 2.} For any probability space and any invertible bimeasurable measure-preserving mixing $T$, OpenAI claims for every fixed $f_1,\ldots,f_4\in L^\infty(\mu)$ that
\[
\frac1N\sum_{k=1}^N\prod_{j=1}^4 f_j(T^{jk}x)\longrightarrow\prod_{j=1}^4\int_X f_j\,d\mu
\quad\mu\text{-almost everywhere}.
\]
Mixing means $\mu(A\cap T^{-m}B)\to\mu(A)\mu(B)$ for every measurable $A,B$ as $|m|\to\infty$; no quantitative rate or standard-space assumption is imposed, and convergence is along all positive $N$.

\paragraph{Mathematical claim 3.} Take an arbitrary probability space and an invertible probability-preserving mixing $T$ with measurable inverse. For each fixed triple of bounded measurable complex functions, OpenAI claims
\[
\frac1N\sum_{k=1}^N f_1(T^kx)f_2(T^{2k}x)f_3(T^{3k}x)
\longrightarrow\prod_{j=1}^3\int_X f_j\,d\mu
\quad\mu\text{-almost everywhere}.
\]
Mixing is the two-set condition $\mu(A\cap T^{-m}B)\to\mu(A)\mu(B)$ as $|m|\to\infty$. All positive $N$ are used and a different triple may have a different excluded null set.

\paragraph{Mathematical claim 4.} For an invertible bimeasurable probability-preserving mixing $T$, any pairwise distinct nonzero integers $\lambda_1,\lambda_2,\lambda_3$ and fixed bounded measurable complex functions $f_j$, OpenAI claims
\[
\frac1N\sum_{k=1}^N\prod_{j=1}^3 f_j(T^{\lambda_j k}x)
\longrightarrow\prod_{j=1}^3\int_X f_j\,d\mu
\quad\mu\text{-almost everywhere}.
\]
Mixing means the two-set correlation condition as $|m|\to\infty$. Negative slopes are allowed, convergence is through all positive $N$, and the excluded null set is allowed to vary when the slopes or functions change.

\subsection{Outline of the argument}
Manuscript 1 (family 154, 2026-10-04): A minimal nonzero centered correlation makes smaller selected spatial lists independent. Conditional-expectation energy saturation stabilizes extensions, and a square joining array is regrouped to create an additional output-only slot; maximality of that slot count gives a contradiction using only the existing at-most-three-slot pointwise theorem. Manuscript 2 (family 154, 2026-10-04): Three-slot oscillation estimates and the companion multiple-mixing theorem supply the analytic inputs. Selected joining laws are packed into iterated arrays, where orthogonality bounds the number of extra-coordinate slots; saturated conditional laws and reanchoring increase that number, contradicting its maximal choice.

Manuscript 3 (family 154, 2026-10-04): The companion trilinear-Hilbert local construction is extended to inputs that change with spatial depth. Its four-term quadratic-phase cancellation yields a block oscillation bound growing like the square root of the block count; failure to be Cauchy would force linear growth, while mixing identifies the pointwise limit. Manuscript 4 (family 154, 2026-10-04): The local progression estimates are extended using an integer left inverse of the complete primitive configuration. Vandermonde cancellation and a height base divisible by every slope difference preserve the Fourier and divisibility bounds through discretization, giving the slope-dependent oscillation estimate needed for convergence.

\subsection{Scope relative to this problem}
Manuscript 1 (family 154, 2026-10-04): Setting unused functions to one gives double convergence when the two commuting maps are distinct powers of a common mixing transformation. Arbitrary commuting transformations without that hypothesis remain outside the theorem.

Manuscript 2 (family 154, 2026-10-04): Setting unused functions to one gives double convergence when the two commuting maps are distinct powers of a common mixing transformation. Arbitrary commuting transformations without that hypothesis remain outside the theorem.

Manuscript 3 (family 154, 2026-10-04): Setting unused functions to one gives double convergence when the two commuting maps are distinct powers of a common mixing transformation. Arbitrary commuting transformations without that hypothesis remain outside the theorem.

Manuscript 4 (family 154, 2026-10-04): Setting unused functions to one gives double convergence when the two commuting maps are distinct powers of a common mixing transformation. Arbitrary commuting transformations without that hypothesis remain outside the theorem.

\subsection{Formal verification}
Family 154 has no Lean scope page in this snapshot; no formalization of these manuscript results is listed.

This site has not executed the Lean files or checked their dependencies. These descriptions report the scope of the linked declarations, rather than a new verification of the proofs.

\subsection{Status caveat}
OpenAI attributes these results to an internal model. Its release notice cautions that results lacking full formalization can contain mistakes. This site has not independently assessed the proofs; a solved label records the scope of the reported claim.

\subsection{Sources}
\begin{itemize}
\item OpenAI, \emph{Pointwise Multiple Ergodic Averages for Mixing Transformations}, 2026-10-04; citation key \texttt{OAI:Pointwise-Multiple-Ergodic-Averages-for-Mixing-Transformations-October-4-2026}. \href{https://github.com/openai/math/blob/adc7f1241b42e322a6451854ab7e4b4c146bf78a/preprints/Pointwise-Multiple-Ergodic-Averages-for-Mixing-Transformations-October-4-2026/multiple-ergodic-averages.pdf}{PDF (GitHub)}; \href{https://github.com/openai/math/raw/adc7f1241b42e322a6451854ab7e4b4c146bf78a/preprints/Pointwise-Multiple-Ergodic-Averages-for-Mixing-Transformations-October-4-2026/multiple-ergodic-averages.pdf}{PDF (direct)}; \href{https://github.com/openai/math/blob/adc7f1241b42e322a6451854ab7e4b4c146bf78a/preprints/Pointwise-Multiple-Ergodic-Averages-for-Mixing-Transformations-October-4-2026/README.md}{Manuscript page and citation}; \href{https://github.com/openai/math/blob/main/preprints/Pointwise-Multiple-Ergodic-Averages-for-Mixing-Transformations-October-4-2026/multiple-ergodic-averages.pdf}{Latest version}; Statement source: \href{https://github.com/openai/math/blob/adc7f1241b42e322a6451854ab7e4b4c146bf78a/preprints/Pointwise-Multiple-Ergodic-Averages-for-Mixing-Transformations-October-4-2026/build/sections/introduction.tex\#L20}{thm:main}.
\item OpenAI, \emph{Pointwise convergence of fourfold ergodic averages for mixing transformations}, 2026-10-04; citation key \texttt{OAI:Pointwise-convergence-of-fourfold-ergodic-averages-for-mixing-transformations-October-4-2026}. \href{https://github.com/openai/math/blob/adc7f1241b42e322a6451854ab7e4b4c146bf78a/preprints/Pointwise-convergence-of-fourfold-ergodic-averages-for-mixing-transformations-October-4-2026/fourfold-ergodic-averages.pdf}{PDF (GitHub)}; \href{https://github.com/openai/math/raw/adc7f1241b42e322a6451854ab7e4b4c146bf78a/preprints/Pointwise-convergence-of-fourfold-ergodic-averages-for-mixing-transformations-October-4-2026/fourfold-ergodic-averages.pdf}{PDF (direct)}; \href{https://github.com/openai/math/blob/adc7f1241b42e322a6451854ab7e4b4c146bf78a/preprints/Pointwise-convergence-of-fourfold-ergodic-averages-for-mixing-transformations-October-4-2026/README.md}{Manuscript page and citation}; \href{https://github.com/openai/math/blob/main/preprints/Pointwise-convergence-of-fourfold-ergodic-averages-for-mixing-transformations-October-4-2026/fourfold-ergodic-averages.pdf}{Latest version}; Statement source: \href{https://github.com/openai/math/blob/adc7f1241b42e322a6451854ab7e4b4c146bf78a/preprints/Pointwise-convergence-of-fourfold-ergodic-averages-for-mixing-transformations-October-4-2026/build/sections/01-introduction.tex\#L21}{thm:main}.
\item OpenAI, \emph{Pointwise convergence of triple ergodic averages for mixing transformations}, 2026-10-04; citation key \texttt{OAI:Pointwise-convergence-of-triple-ergodic-averages-for-mixing-transformations-October-4-2026}. \href{https://github.com/openai/math/blob/adc7f1241b42e322a6451854ab7e4b4c146bf78a/preprints/Pointwise-convergence-of-triple-ergodic-averages-for-mixing-transformations-October-4-2026/pointwise-triple-ergodic-averages-mixing-transformations.pdf}{PDF (GitHub)}; \href{https://github.com/openai/math/raw/adc7f1241b42e322a6451854ab7e4b4c146bf78a/preprints/Pointwise-convergence-of-triple-ergodic-averages-for-mixing-transformations-October-4-2026/pointwise-triple-ergodic-averages-mixing-transformations.pdf}{PDF (direct)}; \href{https://github.com/openai/math/blob/adc7f1241b42e322a6451854ab7e4b4c146bf78a/preprints/Pointwise-convergence-of-triple-ergodic-averages-for-mixing-transformations-October-4-2026/README.md}{Manuscript page and citation}; \href{https://github.com/openai/math/blob/main/preprints/Pointwise-convergence-of-triple-ergodic-averages-for-mixing-transformations-October-4-2026/pointwise-triple-ergodic-averages-mixing-transformations.pdf}{Latest version}; Statement source: \href{https://github.com/openai/math/blob/adc7f1241b42e322a6451854ab7e4b4c146bf78a/preprints/Pointwise-convergence-of-triple-ergodic-averages-for-mixing-transformations-October-4-2026/build/paper/sections/introduction.tex\#L21}{thm:main}.
\item OpenAI, \emph{Triple ergodic averages with distinct integer slopes}, 2026-10-04; citation key \texttt{OAI:Triple-ergodic-averages-with-distinct-integer-slopes-October-4-2026}. \href{https://github.com/openai/math/blob/adc7f1241b42e322a6451854ab7e4b4c146bf78a/preprints/Triple-ergodic-averages-with-distinct-integer-slopes-October-4-2026/triple-ergodic-distinct-slopes.pdf}{PDF (GitHub)}; \href{https://github.com/openai/math/raw/adc7f1241b42e322a6451854ab7e4b4c146bf78a/preprints/Triple-ergodic-averages-with-distinct-integer-slopes-October-4-2026/triple-ergodic-distinct-slopes.pdf}{PDF (direct)}; \href{https://github.com/openai/math/blob/adc7f1241b42e322a6451854ab7e4b4c146bf78a/preprints/Triple-ergodic-averages-with-distinct-integer-slopes-October-4-2026/README.md}{Manuscript page and citation}; \href{https://github.com/openai/math/blob/main/preprints/Triple-ergodic-averages-with-distinct-integer-slopes-October-4-2026/triple-ergodic-distinct-slopes.pdf}{Latest version}; Statement source: \href{https://github.com/openai/math/blob/adc7f1241b42e322a6451854ab7e4b4c146bf78a/preprints/Triple-ergodic-averages-with-distinct-integer-slopes-October-4-2026/build/sections/introduction.tex\#L19}{thm:main}.
\item \href{https://github.com/openai/math/blob/adc7f1241b42e322a6451854ab7e4b4c146bf78a/CONTENTS.md}{OpenAI catalogue entry}.
\end{itemize}
\end{document}
